\section {Potential Architectures} 
\label {sec:approaches}
Our selection of a neural network architecture for FER is guided by several core criteria. Most importantly, the model must achieve high accuracy on low-resolution 64×64 pixel RGB and grayscale images while remaining robust to in-the-wild variations such as pose, illumination, and expression intensity. Given the limited size of typical FER datasets, data efficiency and suitability for training from scratch are should also be considered. 
In addition, model explainability plays a central role. The chosen architecture must be compatible with saliency-based visualization methods in order to highlight facial regions relevant for emotion classification, such as the mouth, eyes, and eyebrows. Finally, while not critical for the initial evaluation, low latency and computational efficiency are considered to support a potential real-time webcam demonstration. Overall, the target architecture represents a balance between accuracy, interpretability, and computational efficiency.

Based on these criteria, we evaluate three representative architecture classes: a convolutional neural network baseline, a transformer-based model, and a lightweight hybrid architecture. 

\subsection{Convolutional Neural Networks}
As a strong and well-established baseline, we select ResNet-18 (or potentially ResNet-50, based on computational resources), which offers a favorable trade-off between performance, model complexity, and interpretability \cite{deepres}. Despite its age, previous research and extensive testing allows for flexible experimentation while still producing respectable results on standard FER datasets. We also plan to explore and implement variants that incorporate attention, such as Squeeze-and-Excitation, Triplet Squeeze-and-Excitation, or Convolutional Block Attention Modules \cite{hu2019squeezeandexcitationnetworks,alhazmi2025achieving3dattentiontriplet,woo2018cbamconvolutionalblockattention}. These changes have been shown to improve the accuracy on FER datasets and can allow the model to better focus on relevant facial regions.

To explore more recent convolutional design principles, we also consider ConvNeXt-Tiny. This architecture adopts transformer-inspired design choices - such as larger kernels, inverted bottlenecks, LayerNorm, and GeLU activations - while retaining convolutional inductive biases, resulting in strong performance on visual recognition tasks \cite{liu2022convnet2020s}. Variants of ConvNeXt have been shown to outperform ResNet architectures on FER benchmarks, while remaining more efficient than full Vision Transformers. Although originally designed for higher-resolution inputs and incurring higher computational cost than shallow CNNs, ConvNeXt-Tiny can be adapted to the 64×64 setting and represents a high-accuracy model with manageable complexity.

Finally, for better implementation of a live webcam emotion recognizer, we also explored more lightweight architectures, with our choice being EfficientNetV2-S \cite{tan2021efficientnetv2smallermodelsfaster}. Through training-aware neural architecture search and improved model scaling, this network reaches competitive accuracy on standard datasets while having much fewer parameters than its competitors. This makes it suitable to train and use for live emotion recognition using a webcam.  

All three of our proposed architectures are convolutional networks with simple feed-forward structures, making them compatible with gradient-based explainability techniques such as Grad-CAM and saliency maps \cite{Selvaraju_2019}.

\subsection{Transformers}
While standard ViTs \cite{dosovitskiy2021vit} struggle with limited data, the \textit{Data-efficient Image Transformer (DeiT)} employs a teacher–student training strategy to improve data efficiency. This allows it to achieve high performance on smaller datasets without needing the massive amount of data and computing power usually required by transformers \cite{Touvron2021DeiT}. Visualization is straightforward to implement, as we can directly extract global attention maps to see the model's focus. To address the loss of local spatial information, we are also considering the \textit{Swin Transformer}. Its hierarchical approach calculates attention within small, shifted windows. This limits the focus to local areas, allowing the model to capture fine-grained details efficiently \cite{Liu2021swin}. This mechanism reintroduces a strong bias for locality which crucial for detecting subtle FER cues like mouth and eye movements while maintaining the global modeling capabilities of transformers. However, creating visual maps is trickier here since local attention maps have to be stitched together to get a complete picture.

\subsection{Hybrid Architectures}
To counter the complexity often introduced by hybrid models, we prioritize lightweight variants designed for efficiency. Most notably, the \textit{Compact Convolutional Transformer} (CCT) is specifically engineered for small-scale datasets and lower resolutions. CCT processes the raw image with convolutions first, which helps it capture important details before passing them to the transformer \cite{Hassani2021CCT}. This architecture is particularly promising, as it provides a solution that is robust, easy to train, and sufficiently low-latency for our real-time webcam demonstration. Visualization is also easy to implement, allowing us to use Grad-CAM for local features and Attention Rollout for high-level emotional cues.